% !TEX program = xelatex
% ============================================================
%  فصل كتاب — Book Chapter Template
%  نموذج فصل كتاب علمي بأسلوب article مع تنسيق الفصل ورأس "Chapter X"
%  Compile with:
%    xelatex main.tex
%    bibtex main
%    xelatex main.tex
%    xelatex main.tex
% ============================================================
%  Features:
%    - Article class styled as a book chapter
%    - Chapter title, author, abstract
%    - Sections, figures/tables placeholders
%    - BibTeX bibliography (apalike style)
%    - Header "Chapter X" via fancyhdr
% ============================================================

\documentclass[11pt,a4paper]{article}

% ---- Geometry ----
\usepackage[a4paper,margin=2.5cm,top=2.5cm,bottom=2.5cm,headheight=16pt]{geometry}

% ---- Core packages ----
\usepackage{graphicx}
\usepackage{array}
\usepackage{booktabs}
\usepackage{tabularx}
\usepackage{multirow}
\usepackage{enumitem}
\usepackage{float}
\usepackage{caption}
\usepackage{titlesec}
\usepackage{fancyhdr}
\usepackage{setspace}
\usepackage{xcolor}
\usepackage{amsmath}

% ---- RTL / Arabic language support (load before hyperref) ----
\usepackage{bidi}
\usepackage[no-math]{fontspec}
\usepackage[quiet]{polyglossia}

% ---- Fonts ----
\setmainfont[Scale=1]{Amiri-Regular.ttf}
\newfontfamily\arabicfont[Script=Arabic,Scale=0.95]{Amiri-Regular.ttf}
\newfontfamily\englishfont[Scale=0.95]{TeX Gyre Termes}
\setmonofont{Amiri-Regular.ttf}
\newfontfamily\arabicfontsf{Amiri-Regular.ttf}

% ---- Languages ----
\setdefaultlanguage[calendar=gregorian,hijricorrection=1]{arabic}
\setotherlanguage[variant=british]{english}

% ---- Hyperref ----
\usepackage[breaklinks,hidelinks]{hyperref}
\hypersetup{
  pdftitle={فصل كتاب},
  pdfauthor={اسم المؤلف}
}

% ---- Captions in Arabic ----
\addto\captionsarabic{%
  \renewcommand{\contentsname}{الفهرس}%
  \renewcommand{\figurename}{الشكل}%
  \renewcommand{\tablename}{الجدول}%
  \renewcommand{\refname}{المراجع}%
  \renewcommand{\abstractname}{الملخص}%
}

% ---- Section formatting ----
\titleformat{\section}{\Large\bfseries}{\thesection}{0.7em}{}
\titleformat{\subsection}{\large\bfseries}{\thesubsection}{0.6em}{}
\titlespacing*{\section}{0pt}{1.4ex plus 0.4ex}{0.9ex plus 0.3ex}
\titlespacing*{\subsection}{0pt}{1.0ex plus 0.3ex}{0.7ex plus 0.2ex}

% ---- Chapter number (manual, for header) ----
% TODO: حدد رقم الفصل هنا.
\newcommand{\chapternumber}{1}
\newcommand{\chaptertitle}{عنوان الفصل (Chapter Title)}

% ---- Headers / footers ----
\pagestyle{fancy}
\fancyhf{}
% TODO: يظهر "الفصل X" في رأس كل صفحة.
\fancyhead[R]{\small\itshape الفصل \chapternumber\ --- \chaptertitle}
\fancyhead[L]{\small اسم المؤلف}
\fancyfoot[C]{\small\thepage}
\renewcommand{\headrulewidth}{0.3pt}
\renewcommand{\footrulewidth}{0pt}

% ---- List spacing ----
\setlist[itemize]{topsep=0.3em, itemsep=0.15em, parsep=0pt}
\setlist[enumerate]{topsep=0.3em, itemsep=0.15em, parsep=0pt}

% ---- Table column types ----
\newcolumntype{L}[1]{>{\raggedright\arraybackslash}p{#1}}
\newcolumntype{C}[1]{>{\centering\arraybackslash}p{#1}}

% ---- Line spacing ----
\setstretch{1.4}

\begin{document}

% ============================================================
% CHAPTER HEADER
% ============================================================
\thispagestyle{fancy}

\begin{center}
{\Large\bfseries الفصل \chapternumber}\\[0.8em]
{\LARGE\bfseries \chaptertitle}\\[1em]
% TODO: اكتب اسم مؤلف الفصل وانتماءه هنا.
{\large اسم المؤلف}\\[0.3em]
{\small القسم، الجامعة}
\end{center}

\vspace{1em}

% ============================================================
% ABSTRACT
% ============================================================
\begin{center}
\begin{minipage}{0.9\textwidth}
\noindent\textbf{الملخص (Abstract):}
% TODO: اكتب ملخص الفصل هنا.

\vspace{0.5em}
\noindent\textbf{الكلمات المفتاحية (Keywords):} كلمة مفتاحية 1، كلمة مفتاحية 2، كلمة مفتاحية 3.
\end{minipage}
\end{center}

\vspace{1em}

% ============================================================
% BODY
% ============================================================

\section{المقدمة (Introduction)}
\label{sec:intro}

% TODO: اكتب مقدمة الفصل هنا. استخدم \cite{bibkey} للإشارة إلى المراجع.
هذا نص تجريبي للمقدمة. يقدم سياق الفصل ضمن الكتاب ويحدد محاوره الرئيسية \cite{example_book}.

\section{الإطار النظري (Theoretical Framework)}
\label{sec:theory}

% TODO: اكتب الإطار النظري هنا.
\subsection{المفهوم الأول}
% TODO: اكتب المحتوى.

\subsection{المفهوم الثاني}
% TODO: اكتب المحتوى.

\section{المنهجية (Methodology)}
\label{sec:method}

% TODO: اكتب منهجية الفصل هنا.

\section{النتائج والتحليل (Results and Analysis)}
\label{sec:results}

% TODO: اكتب النتائج والتحليل هنا.

% ---- Figure placeholder ----
\begin{figure}[H]
\centering
% TODO: استبدل بملف الشكل الفعلي.
\fbox{\parbox[c][3.5cm][c]{0.7\textwidth}{\centering مكان الشكل (Figure placeholder)}}
\caption{عنوان الشكل (Figure caption)}
\end{figure}

% ---- Table placeholder ----
\begin{table}[H]
\centering
\caption{عنوان الجدول (Table caption)}
\begin{tabularx}{\textwidth}{L{4cm} X C{3cm}}
\toprule
\textbf{البند} & \textbf{الوصف} & \textbf{القيمة} \\
\midrule
% TODO: املأ بيانات الجدول.
البند الأول & وصف البند & \LR{100} \\
البند الثاني & وصف البند & \LR{200} \\
\bottomrule
\end{tabularx}
\end{table}

\section{المناقشة (Discussion)}
\label{sec:discussion}

% TODO: اكتب مناقشة النتائج هنا.

\section{الاستنتاجات (Conclusions)}
\label{sec:conclusion}

% TODO: اكتب استنتاجات الفصل هنا.

% ============================================================
% BIBLIOGRAPHY
% ============================================================
% Add your .bib file as bibliography.bib in the project folder.
% Run: xelatex -> bibtex -> xelatex -> xelatex
\bibliography{bibliography}
\bibliographystyle{apalike}

\end{document}
